\section{Los retos de la operación en producción}

\subsection{Por qué DevOps es el «70\% más difícil» de la ingeniería de software}

\begin{importantbox}{La codificación solo ocupa el 30\% del tiempo de ingeniería}
La codificación ocupa apenas alrededor del 30\% del tiempo total de los ingenieros. El 70\% más difícil consiste en ejecutar el código en producción: ahí chocan de lleno la complejidad, los silos de herramientas (tool silos), las brechas de conocimiento y las interdependencias entre los sistemas.
\end{importantbox}

Monitorear y mantener los sistemas en producción sigue siendo una tarea crítica para la misión (mission-critical), que suele recaer en los SRE(Site Reliability Engineer).

\subsection{Las responsabilidades tradicionales de los SRE}

\begin{itemize}
    \item Monitoreo operativo: turnos de guardia (on-call), resolución de fallas (troubleshooting), gestión de infraestructura y seguridad
    \item Resolver un incidente exige reunir información dispersa entre varias fuentes y equipos distintos
    \item Mantener runbooks(manuales operativos)que con frecuencia están desactualizados
    \item Las arquitecturas nativas de la nube (contenerización + Kubernetes) introducen más datos, dependencias y complejidad entre sistemas
\end{itemize}

\begin{warningbox}{El agotamiento de los SRE}
Los SRE sufren agotamiento (burnout) con frecuencia por las rotaciones constantes de turnos de guardia. Que PagerDuty los despierte a las 3 a.m. para atender un repunte de errores 500 es la rutina diaria de muchos SRE. Uno de los valores importantes de los AI SRE es aliviar esta carga humana.
\end{warningbox}

\subsection{Resumen de la sección}

La operación en producción es la parte más demandante en tiempo y más estresante de la ingeniería de software. Los SRE enfrentan múltiples retos, como la fragmentación de la información, la documentación obsoleta y el agotamiento, lo cual abre un enorme espacio para la intervención de la IA.

